\documentclass[11pt,a4paper]{article}
\usepackage[T1]{fontenc}
\usepackage[utf8]{inputenc}
\usepackage{lmodern}
\usepackage[margin=24mm,headheight=15pt]{geometry}
\usepackage{amsmath,amssymb,amsthm,mathtools}
\usepackage{booktabs,longtable,array,tabularx}
\usepackage{xcolor,microtype,enumitem}
\usepackage{fancyhdr}
\usepackage{listings}
\usepackage{xurl}
\usepackage[colorlinks=true,linkcolor=blue!45!black,urlcolor=blue!45!black,citecolor=blue!45!black]{hyperref}
\hypersetup{pdftitle={Porting ProveIt reductions to K230},pdfauthor={Research contribution prepared for ProveIt}}
\setlength{\parskip}{5pt}
\setlength{\parindent}{0pt}
\setlength{\emergencystretch}{2em}
\setcounter{tocdepth}{2}
\pagestyle{fancy}
\fancyhf{}
\fancyhead[L]{\small ProveIt / K230 reduction study}
\fancyhead[R]{\small 19 October 2026}
\fancyfoot[C]{\thepage}
\newtheorem{theorem}{Theorem}
\newtheorem{lemma}[theorem]{Lemma}
\newcommand{\Con}{\operatorname{Con}}
\newcommand{\Halt}{\operatorname{Halt}}
\newcommand{\ZFC}{\mathsf{ZFC}}
\newcommand{\code}[1]{\texttt{\detokenize{#1}}}
\newcommand{\artifact}[1]{\texttt{\detokenize{#1}}}
\lstset{basicstyle=\small\ttfamily,columns=fullflexible,keepspaces=true,breaklines=true,frame=single,rulecolor=\color{black!20}}
\input{results.tex}

\begin{document}
\begin{titlepage}
\vspace*{14mm}
{\Large\sffamily RESEARCH CONTRIBUTION FOR PROVEIT\par}
\vspace{10mm}
{\huge\bfseries Porting ProveIt reductions to K230\par}
\vspace{5mm}
{\LARGE A \CandidateStates-state binary machine calibrated to\par}
{\LARGE $\Con(\ZFC)$\par}
\vspace{12mm}
{\large Exact construction, source equivalence, unbounded runtime\par}
{\large certificates, and alternative encoding experiments\par}
\vspace{12mm}
19 October 2026
\vfill
\begin{abstract}
We transfer compatible reductions from ProveIt's 243-to-235-state line to
Elliot Glazer's independently calibrated 230-state K construction. Two
predecessor-bit reductions are already present in K230. The transfer-plus-one
loop rewrite and its contextual fold give 229 states; delayed erasure gives
a particularly simple 228-state companion, differing from the published
K230 table by two deleted states and eight retained-row edits before renumbering.
A further source-boundary adaptation, using constrained reassociation of a
49-atom closed component and repair of shared decoder contexts, gives 227
states. One more closed source rotation and contextual fold gives the
\textbf{\CandidateStates-state principal candidate}.

The delivery includes exact binary tables, deterministic reconstruction,
independent source and initialization checks, and fresh parametric runtime
proofs. For the principal table, all 6,656 guarded instruction cases close,
using \RuleCount\ symbolic rules and all \TransitionRows\ literal rows.
The exact blank bootstrap reaches the calibrated complete tape after
409,835 transitions. These facts yield an ordinary mathematical proof,
formalizable in ZFC, of nonhalting equivalent to $\Con(\ZFC)$.
The inherited K230 endpoint is checked in Rocq; the reduced table has no
new Rocq or Lean endpoint. A separate equational-query observation produces
a fully serialized 662-bit packet, without asserting a further machine saving.
\end{abstract}
\vfill
\small
Principal upstream: \url{https://github.com/elliotglazer/k230-rocq}\par
Pinned commit: \code{3142e4dd980f0be12f4abb636c4ae2528afc17eb}\par
ProveIt reference delivery: \code{0128d662fd46f72dc84505176ae8ec28affb2d9d}
\end{titlepage}

\begingroup
\small
\setlength{\parskip}{0pt}
\tableofcontents
\endgroup
\clearpage

\section{Result and exact scope}

\begin{theorem}[Ordinary consistency calibration]
Let $M$ be the principal table in \artifact{machines/k226.tm}, with
working states $0,\ldots,225$, initial state $0$, and a separate halt
outcome $226$. Its tape alphabet is $\{0,1\}$, its initial tape is
identically zero on $\mathbb Z$, and every transition moves one cell
left or right. The constructions and proofs in this delivery establish
the ordinary mathematical statement
\[
 \ZFC\ \vdash\
 \bigl(\neg\Halt(M,\mathrm{blank})\ \longleftrightarrow\ \Con(\ZFC)\bigr).
\]
Here $\Con(\ZFC)$ means that no finite contradiction is derivable in
the standard classical first-order ZFC proof system identified by the
inherited K230 calibration.
\end{theorem}

The exact byte identity of the principal canonical table is
\begin{center}\small\ttfamily\CanonicalHash\end{center}
The inherited-label table, more convenient for comparing old local
reductions, has SHA-256
\begin{center}\small\ttfamily\InheritedHash\end{center}
Its start, dispatcher and halt labels are 95, 96 and 254. Those labels
are names: there are exactly \CandidateStates\ working states. In
particular, inherited working labels may exceed 254; 254 alone is the halt.
The state map supplied with each table verifies the bijective renaming.
Initialization is included in the state count, with no supplied program
on the tape and no separate free loader.

The report calls $M$ a candidate to preserve the project's distinction
between an ordinary mathematical construction and a new proof-assistant
endpoint. The evidence is stronger than a bounded execution experiment:
each instruction proof quantifies over arbitrary natural register values
and arbitrary finite spare storage. It is weaker as a mechanically
trusted artifact than compiling a new Rocq theorem for $M$. The latter
has not been done. No explicit internal PA derivation or serialized
ZFC derivation is supplied either.

\subsection{Why the theorem follows}

The argument has four separately checked interfaces:
\begin{enumerate}[leftmargin=*,itemsep=2pt]
\item The pinned K230 formal theorem identifies blank halting with a
finite StandardZFC refutation. The actual literal Rocq register program
is independently matched to the published instruction catalogue.
\item Four local loop rewrites preserve whole-program halting for
every initial natural register file. Subsequent reassociation preserves
the ordered instruction labels and every guarded successor exactly.
\item The principal binary table implements every guarded instruction
of that source on its complete quantified tape frame, with positive
finite duration and no unaccounted interior halt.
\item Its exact blank startup reaches an instance of that frame with
the actual nonzero seed from K230.
\end{enumerate}

Thus, as statements about finite computations and finite proofs,
\[
\begin{aligned}
 \Halt(M,\mathrm{blank})
 &\longleftrightarrow \Halt(P_{\rm new},R_0)\\
 &\longleftrightarrow \Halt(P_{230},R_0)\\
 &\longleftrightarrow \mathsf{StandardZFC}\vdash\bot.
\end{aligned}
\]
Negating gives the desired equivalence. The new transport uses induction
on natural counters, finite source executions, unary words and finite
lists; these are ordinary arguments available in ZFC. The inherited
logical calibration concerns the precise finite-refutation predicate,
not an unspecified semantic notion of consistency.
The ZFC-provability assertion uses the inherited finitary first-order
proof translations together with these new elementary inductions.
For example, \code{KUpperConsistency.v} records finite axiom lists and
actual first-order derivations in \code{core_holdsL}, and its final
consistency transfer assembles a finite contradiction. This is the
object-level mathematical content being reused; the upstream Rocq
kernel check certifies that implementation of the proof.\footnote{
\url{https://github.com/elliotglazer/k230-rocq/blob/3142e4dd980f0be12f4abb636c4ae2528afc17eb/coq/settheory/KUpperConsistency.v}.}

\subsection{What transferred from 243 to 235}

The earlier eight savings are mechanisms with hypotheses, not eight
interchangeable deletions. Table~\ref{tab:transfer} states the outcome
of inspecting their actual source patterns and transition ports.
The predecessor and loop facts are documented in ProveIt's maintained
241- and 240-state pages; the next three folds are documented in its
237-state page.\footnote{
\url{https://github.com/VladimirReshetnikov/ProveIt/blob/aaf47dea43ec6cc3c1f42b6ffd56a92fb6e7fdab/Computability/BusyBeaver/Research/ZFC/candidate-fold241/README.md};
\url{https://github.com/VladimirReshetnikov/ProveIt/blob/aaf47dea43ec6cc3c1f42b6ffd56a92fb6e7fdab/Computability/BusyBeaver/Research/ZFC/candidate-rotation240/README.md};
\url{https://github.com/VladimirReshetnikov/ProveIt/blob/aaf47dea43ec6cc3c1f42b6ffd56a92fb6e7fdab/Computability/BusyBeaver/Research/ZFC/candidate-composed237/README.md}.}

\begin{table}[htbp]
\small
\begin{tabularx}{\textwidth}{@{}l l X@{}}
\toprule
Old step & Mechanism / inherited states & Outcome on K230 \\
\midrule
243 $\to$ 242 & $27\to15$ & Already present; no new saving.\\
242 $\to$ 241 & $39\to17$ & Already present; no new saving.\\
241 $\to$ 240 & Loop rewrite; $90\to72$ & Transfers: K230 $\to$ K229.\\
240 $\to$ 239 & Source boundary; split 218 & Adapts with K-specific source/context repair; one further saving.\\
239 $\to$ 238 & Delayed erasure; $169\to80$ & Transfers; combines with both sources.\\
238 $\to$ 237 & Shared helper; $235\to98$ & Required source/helper pattern is absent in the inspected K230 representation.\\
237 $\to$ 236 & Rotation, then $223\to182$ & Depends on a different source boundary layout; no further accepted saving here.\\
236 $\to$ 235 & Source rotation and context fold & The method adapts with a different rotation and state pair: $508\to86$, giving K227 $\to$ K226.\\
\bottomrule
\end{tabularx}
\caption{Transfer audit. ``Absent'' or ``no accepted saving'' is a statement
about these inspected constructions, not an impossibility result.}
\label{tab:transfer}
\end{table}

Four additional savings therefore survive in the principal construction:
230 to 226. This is nine fewer states than the requested 235-state
reference, and seventeen fewer than the earlier 243-state table. The
unmodified logical theory remains Glazer's K; the contribution here is
the calibrated composition and the boundary adaptation.

\input{k230_foundation.tex}
\input{source229_proof.tex}

\section{A compact and independently useful 228-state companion}
\label{sec:228}

Before the more substantial source-boundary work, the literal loop and
delayed-erasure ports already give a small, easily reviewed result.
The complete table is \artifact{machines/k228.tm}; its inherited
version is \artifact{machines/k228_old.tm}. The canonical SHA-256 is
\begin{center}\small\ttfamily
06726bd731f3a42e3ccd1c4189d80af74ba2d486e333045097480fcaa3313001.
\end{center}
It is retained as a companion because its patch is much smaller than
the principal source-recompiled construction.

\subsection{Eight retained rows in the original K230 numbering}

Use the published K230 labels in this subsection, before any renumbering.
Delete states 79 and 140. Replace exactly these eight retained rows:
\begin{center}
\begin{tabular}{rrrrr}
\toprule
State & Read & Write & Move & Next\\
\midrule
121&0&0&R&63\\
121&1&1&R&13\\
63&0&0&L&123\\
63&1&1&L&30\\
70&1&1&R&63\\
71&0&0&L&91\\
80&1&1&L&138\\
126&0&0&R&71\\
\bottomrule
\end{tabular}
\end{center}
Every other retained row stays unchanged. Reindex the remaining working
states in their original order as $0,\ldots,227$, and change the old
halt label 230 to 228. The supplied JSON patch and state map make this
description executable and check its final byte digest. A source rewrite
is still essential to its proof: these eight rows are not justified
merely by visually comparing the machine graphs.

\subsection{The contextual fold from 90 into 72}

We return to inherited labels. The source change alters only
$142/0:0R37\to0R72$ and $142/1:1R72\to1R17$ in the 230-state
intermediate. This changes the context in which 72 is entered.
The independently recomputed invariant gives a 1 immediately left of
55, 142 and 90, and the bits 10 immediately left of 72.
The state map sends 90 to 72 and fixes every other state. Replace the
rows of 72 by $0L144$ and $1L36$, redirect $79/1$ to 72, and delete 90.

For every admitted initial tape and head position, each old transition
is reproduced by a positive new block with the same whole tape and
mapped endpoint. There are 456 one-step cases, three three-step cases,
and one nine-step case, covering all 460 old rows. The checker explores
934 symbolic configurations and rederives the predecessor invariant
from the literal table. Its four exceptional blocks are:
\begin{center}\small
\begin{tabular}{lll}
\toprule
Old row & New control path & Length\\
\midrule
72/0 & $72,144,47,15,32,46,17,106,93,57$ & 9\\
72/1 & $72,36,53,64$ & 3\\
90/0 & $72,144,82,93$ & 3\\
90/1 & $72,36,54,71$ & 3\\
\bottomrule
\end{tabular}
\end{center}
The path includes both initial and final controls. The proof does not
restrict the initial tape to blank; the source-changing step before
this fold is separately calibrated for the intended blank initialization.

\subsection{Delayed erasure and its precise proof premise}

Delete inherited state 169 into 80. Change $80/0$ to $0L106$,
$91/1$ to $1L167$, and redirect $147/0$ to $0R80$. The 1 at the
second changed row is retained temporarily; existing states can use it
to distinguish two control roles before restoring the required tape.

The independent local checker proves these complete tape identities:
\begin{center}\small
\begin{tabularx}{\textwidth}{@{}l X r@{}}
\toprule
Old row & Additional local premise & New transitions\\
\midrule
80/0 & Left neighbor is 1 & 5\\
169/0 & Left neighbor is 0 & 3\\
169/1 & Left neighbor is 0 and right neighbor is 1 & 5\\
91/1 & The preceding three cells, left to right, are 011 & 9\\
\bottomrule
\end{tabularx}
\end{center}
The remaining 454 rows commute in one step. The left-neighbor premises
follow from a syntactic inductive invariant. The right-neighbor premise
for $169/1$ does not follow from that invariant alone; an all-tape
claim based only on those predecessor facts would be unjustified.

For the blank prefix, all six relevant $169/1$ occurrences have the
required right-hand 1. The independent checker compares all 409,692
K229 startup checkpoints to their variable-duration K228 images.
The four exceptional counts are respectively 18, 12, 6 and 3, giving
\[
409691 + 18\cdot4 + 12\cdot2 + 6\cdot4 + 3\cdot8=409835.
\]
For the unbounded part, the accepted argument proves every instruction
contract directly on the final K228 table. This eliminates the need
to assume the extra guard for arbitrary reachable tapes. All 6,656
symbolic cases pass, with 1,614,426 rules and all 456 rows used.

\input{boundary227.tex}
\input{extension226.tex}

\section{Unbounded binary runtime proof}
\label{sec:runtime}

\subsection{Exact quantified tape frames}

For a natural $v$, put $B(v)=0\,1^{v+1}$ and $C(v)=1\,0\,1^v$;
extend both encodings to finite lists by concatenation. At inherited
dispatcher 96 the head is immediately right of a left marker 1, the
exterior to its left is zero, and the tape to its right is
\begin{equation}
 p\;0^{64+g-|p|}\;1\;
 B(V_0)\cdots B(V_8)\;B(s)\;B(A)\;0^\omega.
 \label{eq:frame}
\end{equation}
Here $p$ is a literal source-leaf path, $g,s\in\mathbb N$, and $A$
is any finite list of naturals, including the empty list. The nine
logical registers are independent, with physical permutation
\[
 V_{\pi(i)}=R_i,\qquad \pi=(6,2,0,4,3,7,8,5,1).
\]
The explicit $B(s)$ supplies the nonempty spare list needed by this
interpreter. The accepted source has maximum depth \SourceDepth,
so the gap exponent is nonnegative for every allowed $g$.

An increment raises its logical register by one and appends one
$B(1)$ spare block. A positive decrement lowers its register by one
and adds one to the gap parameter. Other operations preserve registers
and spare blocks. A program-counter overflow moves the marker and
head one cell left and adds one more to the gap parameter. Thus the
frame family remains closed: the gap never becomes negative and the
spare list never becomes empty.

For each decrement, the zero case fixes the selected register at zero
and the positive case writes it as $1+n$ for an arbitrary natural $n$.
The other eight registers are still independent. All other instructions
require one case. This is the source of the exact 6,656-case coverage.

\subsection{Sound rules of the executable checker}

The checker represents both sides of the head by runs and opaque
finite-list atoms. Every run length is an affine expression in thirteen
slots: constant, extra gap, nine primary parameters, first spare, and
remaining-list width. It requires nonnegative coefficients. A literal
step from a run not handled by a scan loop requires constant coefficient
at least one, so the read cell exists for every valuation.

There are three execution rules. A literal transition is just the
corresponding row of the supplied table. A homogeneous self-loop scan
is justified by induction on its arbitrary natural run length,
including the zero-length case. A traversal of an opaque finite list
is justified by structural induction after checking a one-block lemma
on the actual table, with arbitrary opaque contexts on both sides.
The four one-block lemmas are:
\begin{center}
\begin{tabular}{rlll}
\toprule
Entry & Direction & Encoding before $\to$ after & Duration\\
\midrule
73 & Right & $B(v)\to B(v)$ & $v+2$\\
95 & Right & $B(v)\to C(v)$ & $v+2$\\
48 & Left & $B(v)\to B(v)$ & $v+2$\\
112 & Left & $C(v)\to B(v)$ & $v+2$\\
\bottomrule
\end{tabular}
\end{center}
Each lemma returns to the same entry state beyond the entire block
and leaves both exterior contexts untouched. List induction therefore
covers every finite tail. Its width parameter is instantiated as
$W(A)=\sum_{v\in A}(v+2)$; treating this nonnegative number symbolically
does not exclude any actual finite list.

Normalization uses exact word identities: merge adjacent equal runs,
remove zero-length runs, absorb finite trailing zeros into a blank
exterior, and use
\[
 C(A)1=1B(A),\qquad
 1\operatorname{rev}(C(A))=\operatorname{rev}(B(A))1.
\]
The first follows from $C(v)1=1B(v)$ by induction on $A$; the second
is its reversal. No unknown nonzero tape data are discarded.

The generic engine and helper are retained byte for byte from the
ProveIt evidence. Their original command-line entry point pins an older
source. The new small adapter validates the K source and calls the
unchanged generic \code{prove_blocks} and \code{case} functions.
It checks the actual complete binary table, source prefix tree,
register bounds, source depth, control operands and all program-counter
normalizations before proving every case. The adapter does not invoke
the old CLI with its unrelated source hash removed.

\subsection{Program-counter normalization}

A primitive source step can leave a path shorter or longer than the
next leaf path. Follow its bits, then zeros, down the complete source
tree. The independent source check verifies, for every possible
guarded successor, that any discarded suffix contains only zeros.
Appending or removing those zeros changes only the division between
$p$ and its zero padding in \eqref{eq:frame}; it leaves every tape cell
and the header position unchanged. This finite check supplies the
``no stray bit'' interface needed to concatenate the per-leaf proofs.
The initial empty path resolves in the same way to the first leaf,
\code{000}.

\subsection{Why a successful case proves actual progress}

For a nonhalting instruction, a case closes only at the exact complete
target frame, including its whole tape, state, relative head position,
registers, arbitrary existing spare list and any appended spare.
The accumulated absolute displacement is checked separately. A
premature halt fails a nonhalting case. For the unique source halt,
reaching the separate halt outcome is the required result; the final
tape contents are irrelevant to halting.

The duration is an affine expression with nonnegative coefficients
and a positive constant. The smallest certified constant is 13.
Each instruction therefore takes at least 13 concrete transitions,
for every allowed valuation. If the source runs forever, these
positive durations concatenate to an infinite machine run. If the
source halts, the finite sequence of positive blocks reaches the halt
outcome. Conversely, determinism and the absence of interior halt
ensure that a finite machine halting run comes from a finite source
execution ending in the halt leaf. There can be no infinite sequence
of simulated zero-time instructions before a finite machine halt.

The engine's 4,000-iteration cap bounds construction of a symbolic
derivation. It does not bound the concrete execution, register values,
gap, number of spare blocks or their sizes. In the principal proof
the longest symbolic case uses only \MaximumSymbolic\ rules.

\subsection{The actual blank seed}

Two independently implemented tape representations---a sparse set of
one-cell positions and a two-list zipper---compute the complete blank
bootstrap. After 409,835 transitions they agree exactly on state 96,
head $-5$, marker $-6$, header 176, and the entire support of 351 one
cells. Every other cell is zero. The endpoint is reconstructed from
the literal K230 bootstrap recipe, not inferred from a bounded
prefix of the tape.

It is exactly \eqref{eq:frame} with
\[
 R_0=(3,0,0,0,3,0,0,0,0),\qquad
 g=117,\quad s=1,\quad A=[1]^{166}.
\]
Hence the header distance is $64+117=181$, the physical primary
values are $(0,0,0,3,0,0,3,0,0)$, and there are 167 spare blocks
of value 1. The initial empty program counter and its first resolved
leaf \code{000} denote the same tape. Earlier visits to state 96
during initialization need not be clean dispatcher frames; the proof
uses the specified complete endpoint.

\section{Validation, reproduction and further search}

\begin{table}[htbp]
\centering\small
\begin{tabular}{@{}lrrr@{}}
\toprule
Verified property & Simple K228 & K227 & Principal K226\\
\midrule
Working states & 228 & 227 & 226\\
Literal rows & 456 & 454 & 452\\
Source leaves & 5,121 & 5,121 & 5,121\\
Guarded cases & 6,656 & 6,656 & 6,656\\
Distinct source subtrees & 212 & 211 & 211\\
Maximum source depth & 40 & 37 & 37\\
Maximum zero run & 6 & 6 & 6\\
Symbolic rule applications & 1,614,426 & 1,591,422 & \RuleCount\\
Minimum concrete instruction duration & 13 & 13 & 13\\
Maximum symbolic case length & 415 & 388 & \MaximumSymbolic\\
Exact blank bootstrap & 409,835 & 409,835 & 409,835\\
Final tape one-cells & 351 & 351 & 351\\
\bottomrule
\end{tabular}
\caption{Fresh complete verification, with natural-valued parameters and
finite lists quantified by induction. Row coverage alone is not the proof.}
\end{table}

\subsection{Portable replay}

From the extracted archive root, Python 3.10 or later suffices:
\begin{lstlisting}
python3 -B reproduce.py --output-dir /tmp/k226-replay
\end{lstlisting}
This reconstructs the machines and source, checks the pinned upstream
source and retained CI artifact, repeats the literal Rocq-source bridge,
checks the source transformations, recomputes the bootstrap and local
fold evidence, and reruns the complete principal symbolic proof.
Assertions must remain enabled: do not use \code{-O} or set
\code{PYTHONOPTIMIZE}. No network or external package is needed.

For a quick construction and finite-interface check:
\begin{lstlisting}
python3 -B reproduce.py --quick --output-dir /tmp/k226-quick
\end{lstlisting}
The quick mode also audits the integrity and complete case coverage of
the retained runtime records. It explicitly does not claim a fresh
universal derivation. A \code{PASS} string or a hash match alone is
not a substitute for rerunning the full induction engine. Add
\code{--all-runtime} to the full command to rederive the K228 and K227
companion proofs as well as the principal K226 proof.

The original K230 Rocq snapshot and its attribution are retained under
\artifact{vendor/k230-rocq}. Its upstream mathematical dependencies
are pinned by its own manifest. A new local Rocq rebuild requires the
toolchain and separately fetched dependencies described in that
snapshot. The ordinary Python replay neither runs nor replaces that
Rocq build.

\subsection{Rejected attempts and the useful next direction}

Several initially attractive source rotations failed a concrete
instruction contract. A pure decoder state can also be entered while
moving left as an arithmetic helper. Reparenthesizing it without
preserving that context can return to the wrong source successor even
when the abstract register-program graph is unchanged. In one rejected
layout the first failing leaf was zero-guard decrement of register 1
at ordinal 568. A later candidate exposed a missing helper for state
218. These were genuine local counterexamples to those compiler
layouts; they were not evidence that the source reassociation was
mathematically invalid.

The accepted boundary construction makes incoming tape tags and helper
roles explicit constraints. That is the main reusable methodological
change: optimize the source tree jointly with the implementation's
context requirements. A source DAG with fewer nodes is a promising
proposal, but it is only useful after the complete binary table passes
the quantified runtime contract.

Restricted direct searches on the simple K228 table found no additional
fold at 17- or 33-transition budgets, including selected head-offset
and one-bit tag variants. A separate correlated-window/helper search
at radius 5 and block budgets 65 inspected 25,878 state pairs, with
516,060,285 helper comparisons, and produced no candidate in that
family. These negative searches are retained as discovery records;
their finite limits are not lower bounds on all possible machines.
The successful dynamic program demonstrates why changing the source
and its shared-state contexts can be more productive than extending
only those direct-fold searches.

After accepting K226, a final targeted search examined 15, 120 and 680
distinct layouts at depths one, two and three additional closed
rotations, respectively, keeping every intermediate zero run at most
six and freezing startup rows. No further generic context fold appeared.
An earlier continuation allowed an intermediate seven-zero run and
examined 1,602 layouts through depth two, also without a hit. The
construction is frozen at 226; these searches make no global
minimality claim.

\input{alternatives.tex}

\section{Attribution and integration notes}

The compact theory K, its numerical axiom encoding, the 230-state
table and its formal consistency calibration remain credited to
Elliot Glazer and the K230 development. The predecessor-bit, loop,
boundary and delayed-erasure mechanisms are inherited from the
identified ProveIt constructions and incoming contributions. The new
work consists of applying their actual hypotheses to K230, composing
the compatible savings, finding and checking the K-specific boundary
layout, and supplying the exact reproducible reduced tables.

The calculus and proof-search architecture trace to Mira-Cult; the
Rocq reference development to Jeff Ketchersid's busy-beaver project;
and the compiler framework to Andrew James Wade's Turing Machine
Explorer. BusyCoq, LibTactics and the Coq Library of Undecidability
retain their upstream notices. ProveIt contributors retain their
existing attribution. No personal author is inferred for an anonymous
incoming contribution.

The delivery preserves ProveIt's MIT No Attribution notice and the
K230 \artifact{THIRD_PARTY.md} and \artifact{LICENSES} files. The
latter expressly does not assign a blanket license to custom upstream
or adapted files whose scope was not specified. The new report does
not relicense those files. The snapshot contains source and build
instructions rather than precompiled proof objects.

For incorporation, the principal candidate's table, source, state map,
construction recipe and executable proof records should travel
together. The simple K228 companion is useful for a separate minimal
diff review. The retained upstream CI logs establish exactly which
formal baseline is being used. The final \artifact{MANIFEST.sha256}
identifies every shipped artifact byte; generated replay output belongs
in a separate directory.

\appendix
\input{e662.tex}

\section{Complete canonical transition table}
\label{sec:table}

The following table is generated directly from
\artifact{machines/k226.tm}. An entry $wDq$ means: write $w$, move
one cell in direction $D\in\{L,R\}$, and enter $q$. The symbol $H$
denotes the separate halt outcome 226, which is excluded from the
226 working states. The initial state is 0. Each working state has
exactly one row for each read symbol; no stationary move or larger
alphabet is used.

\input{transition_table.tex}

\end{document}
